\SourceLeaf{59}
\begin{center}
{\large NEUE BEWEISE UND ERWEITERUNGEN\par}
\smallskip DES\par
\smallskip{\large FUNDAMENTALTHEOREMS\par}
\medskip IN DER THEORIE DER QUADRATISCHEN RESTE.\par
\end{center}
\BookRule{20mm}
Das Fundamentaltheorem über quadratische Reste, das zu den schönsten Wahrheiten der höheren Arithmetik gehört, wurde zwar durch Induktion leicht entdeckt, aber mit weit größerer Mühe bewiesen. In diesem Gebiet kommt es häufig vor, dass die Beweise der einfachsten Wahrheiten, die sich dem Forschenden durch Induktion wie von selbst darbieten, tief verborgen liegen und erst nach vielen vergeblichen Versuchen endlich ans Licht gebracht werden können, vielleicht auf einem ganz anderen Weg als dem, auf dem man sie gesucht hatte. Dann geschieht es nicht selten, dass sich, sobald ein Weg gefunden ist, nach und nach \textit{mehrere} weitere eröffnen, die zum selben Ziel führen: einige kürzer und unmittelbarer, andere gleichsam auf Umwegen und ausgehend von ganz anderen Grundsätzen, zwischen denen und der gestellten Frage man kaum einen Zusammenhang vermutet hätte. Dieser erstaunliche Zusammenhang zwischen tiefer verborgenen Wahrheiten verleiht solchen Betrachtungen nicht nur einen besonderen Reiz; er verdient auch deshalb sorgfältige Untersuchung und Aufklärung, weil sich daraus nicht selten neue Hilfsmittel oder Fortschritte für die Wissenschaft selbst ergeben.

Obwohl daher der hier behandelte arithmetische Satz durch die früheren Bemühungen, die vier völlig verschiedene Beweise geliefert haben\begingroup\renewcommand{\thefootnote}{*)}\footnote{Zwei sind im vierten und fünften Abschnitt der \textit{Disquisitiones Arithmeticae} dargelegt, der dritte in einer besonderen Abhandlung (\textit{Commentt. Soc. Gotting. Vol. XVI}); der vierte ist in die Abhandlung \textit{Summatio quarundam serierum singularium} aufgenommen (\textit{Commentt. Recentiores, Vol. I}).}\endgroup, vollständig\SourceLeaf{60} erledigt erscheinen mag, kehre ich dennoch zu demselben Gegenstand zurück und füge zwei weitere Beweise hinzu, die gewiss neues Licht darauf werfen werden. Der erste ist dem dritten insofern verwandt, als er von demselben Lemma ausgeht; danach verfolgt er jedoch einen anderen Weg, so dass er mit Recht als neuer Beweis gelten kann, der an Eleganz jenem dritten wenn nicht überlegen, so doch keineswegs unterlegen erscheinen wird. Der sechste Beweis beruht dagegen auf einem völlig anderen, feineren Grundsatz und liefert ein neues Beispiel für den erstaunlichen Zusammenhang zwischen arithmetischen Wahrheiten, die auf den ersten Blick sehr weit voneinander entfernt liegen. Diesen beiden Beweisen ist ein neuer, sehr einfacher Algorithmus beigegeben, um zu entscheiden, ob eine gegebene ganze Zahl ein quadratischer Rest oder Nichtrest einer gegebenen Primzahl ist.

\newpage
Es gab noch einen anderen Grund, die schon vor neun Jahren versprochenen neuen Beweise gerade jetzt zu veröffentlichen. Als ich im Jahr 1805 begann, die Theorie der kubischen und biquadratischen Reste zu untersuchen, einen weit schwierigeren Gegenstand, erlebte ich nahezu denselben Verlauf wie einst in der Theorie der quadratischen Reste. Die Sätze, die diese Fragen vollständig erschöpfen und eine erstaunliche Analogie zu den Sätzen über quadratische Reste erkennen lassen, wurden sogleich durch Induktion entdeckt, sobald sie auf einem geeigneten Weg gesucht wurden; alle Versuche jedoch, in jeder Hinsicht vollständige Beweise dafür zu gewinnen, blieben lange Zeit vergeblich. Gerade dies spornte mich an, den schon bekannten Beweisen über quadratische Reste mit solchem Eifer immer weitere hinzuzufügen, getragen von der Hoffnung, unter vielen verschiedenen Methoden könne die eine oder andere etwas zur Erhellung des verwandten Gegenstands beitragen. Diese Hoffnung war keineswegs vergeblich, und auf unermüdliche Arbeit folgten schließlich gute Ergebnisse. Bald wird es möglich sein, die Früchte meiner Nachtarbeit ans Licht der Öffentlichkeit zu bringen; bevor ich jedoch dieses schwierige Werk in Angriff nehme, habe ich beschlossen, noch einmal zur Theorie der quadratischen Reste zurückzukehren, alles noch darin zu Erledigende abzuschließen und so von diesem Teil der höheren Arithmetik gleichsam Abschied zu nehmen.
